\begin{textsample}{POS Dim 3 – vem\_ed\_10 – Score 5.00 – vem\_ed\_10/t041.txt}  \label{ex:f3_pos_050}
Jornada de PCIs: internacionalização em casa Após a abertura, seguiram-se os \textbf{painéis} \textbf{temáticos}.
Em “Gestão da Internacionalização em Casa”, Adriane Fontana, diretora da Fatec São Caetano, comentou que a gestão dos PCIs se dá com três pilares: divulgação, motivação e envolvimento de todos os atores envolvidos: direção, coordenação, professores e alunos.
Ruy Accioly, coordenador do curso de Logística da Fatec Baixada Santista, indicou os principais desafios na \textbf{realização} do PCI com a Inacap (Chile), sobre economia circular: idioma, tecnologia da informação e integração das equipes.
Dentre os principais resultados obtidos a partir do PCI, destacou o aprendizado de novas tecnologias da informação, publicação de artigos e TCCs.
Sandro Pinto, coordenador do curso de Gestão da Qualidade da Fatec Lins, enviou um vídeo sobre as fases da gestão dos PCIs: planejamento das atividades, envolvimento dos professores e engajamento dos estudantes.
Regiane Camargo, professora de espanhol, relatou os projetos da Fatec Guaratinguetá: com a Uniminuto (Colômbia), uma \textbf{análise} linguística e semiótica sobre anúncios comerciais em redes sociais; com a Inacap (Chile), uma apresentação em português e espanhol sobre uma empresa; com a Universidade de Aveiro (Portugal), apresentação de um pitch, e um projeto de conversação em língua inglesa com a SUNY Delhi (EUA).

% matched lemmas: análise, painel, realização, temático
\end{textsample}
